\section{Discussion}

The measurements support one claim, and it is narrower than the project's
original question. Whether a learned policy is worth having depends on how
fast the workload moves, and that condition is now measured rather than
assumed. A learned policy is worth building under fast drift and is a waste of
effort under slow drift.

The win over the rule is real but narrow, and it is a win over an opponent
that has itself stopped working. The agent reaches it by being slower and much
cheaper, which is what the combined cost was written to reward, so this is the
objective working as intended rather than a surprise. It is not evidence that
the agent has learned to place data well. The cost-free reference sits at
0.813 against the agent's 1.486, and that distance, not the distance to the
rule, is the measure of what is left.

One methodological point generalises beyond this project. Comparing a learned
policy with a heuristic on latency alone compares them on something only one
of them is optimising. Reporting latency alone would make the agent look worse
than it is, and reporting movement alone would make it look better. Report the
two together, or report the combined cost.
